%
% Documento: Fundamentação teórica
%

\chapter{Fundamentação teórica}
\label{cap:fundamentacao}

\section{RFID aplicado ao inventário patrimonial}

Um sistema RFID é formado, em sua configuração básica, por um interrogador, uma ou mais antenas e etiquetas associadas aos objetos. O interrogador transmite energia e comandos; a tag recebe o sinal, processa a solicitação e devolve seu identificador. Em aplicações patrimoniais, o EPC pode funcionar como chave de associação ao registro administrativo do bem. Os atributos do ativo não precisam ser gravados integralmente na tag, pois podem permanecer em um sistema externo ligado ao identificador.

Etiquetas RFID podem ser ativas, semipassivas ou passivas. As tags ativas possuem fonte própria de energia e podem alcançar distâncias maiores, mas aumentam custo e manutenção. As etiquetas passivas são energizadas pelo campo do leitor, dispensam bateria e são adequadas à identificação em escala \cite{chawla2007passive}. Este trabalho utiliza etiquetas passivas UHF, pois o objetivo é conferir vários bens a uma distância superior à obtida com NFC.

RFID não elimina automaticamente os problemas do inventário. Uma tag fora da zona de leitura, aplicada sobre material inadequado ou orientada em um nulo do campo pode deixar de responder. Também podem ocorrer leituras externas ao ambiente que se pretendia conferir. Portanto, a utilidade do sistema depende tanto da taxa de identificação dos bens presentes quanto do controle espacial da zona de leitura.

\section{NFC, HF e RFID UHF}

NFC opera em 13,56~MHz e utiliza acoplamento de campo próximo. Sua distância reduzida é desejável em pagamentos, autenticação e interações intencionais, mas obriga o operador a aproximar o dispositivo de cada objeto. RFID UHF passivo opera, conforme a região, em porções da faixa entre aproximadamente 860 e 960~MHz e utiliza propagação em campo distante e retroespalhamento, o que permite leituras a alguns metros \cite{turner2007nearfield,rao2005uhf}.

No retroespalhamento, a etiqueta altera a impedância ligada à própria antena e, consequentemente, modula a parcela do sinal refletida ao interrogador. A distância útil depende do enlace de ida, que deve fornecer energia suficiente ao chip, e do enlace de retorno, que deve produzir sinal detectável no receptor. Potência configurada, ganho, perdas em cabos e conectores, sensibilidade da tag, polarização e multipercurso participam desse resultado \cite{wang2009uhf,thomas2012qam}.

Por essa razão, ``alcance'' não é uma constante do leitor. Trata-se de uma propriedade do conjunto formado por leitor, antena e tag em uma configuração e ambiente específicos. A distância máxima de uma leitura isolada também não equivale a alcance operacional; para inventário, interessa a distância na qual a identificação é repetível.

\section{Interface aérea e regulamentação}

\subsection{ISO/IEC 18000-63 e EPC UHF Gen2}

A ISO/IEC 18000-63 especifica a interface aérea para RFID UHF passivo do tipo C, na faixa de 860 a 960~MHz \cite{iso1800063}. O protocolo EPC UHF Gen2 da GS1 é harmonizado com esse padrão e define requisitos físicos, comandos, estados, sessões, seleção de populações e acesso às memórias das tags \cite{gs1gen2}.

A compatibilidade com esse protocolo é o critério principal para interoperabilidade entre o YPD-R200 e a etiqueta AD-239 M730. Este trabalho restringe a análise a essa etiqueta, pois o desempenho continuará dependente do projeto da antena do inlay, da instalação no bem e da orientação em relação à antena APCA8090.

\subsection{Uso do espectro no Brasil}

No Brasil, equipamentos RFID enquadram-se nos requisitos de radiação restrita. O Ato Anatel n.\ 14.448/2017, em sua redação consolidada, apresenta faixas e limites aplicáveis a sistemas RFID, incluindo de 902 a 907,5~MHz e de 915 a 928~MHz para interrogadores UHF \cite{anatelrfid}. A faixa ampla de 860 a 960~MHz impressa em antenas e etiquetas indica cobertura de projeto, mas não autoriza a transmissão em todo esse intervalo.

O módulo deve ser configurado para a região correta e o equipamento utilizado em uma implantação deve ser homologado ou integrar produto cuja certificação seja aceita pela Anatel. Ensaios acadêmicos não justificam operar em canais inadequados ou exceder os limites de emissão. A configuração regional deve ser registrada em cada experimento, assim como potência, antena e ambiente.

\section{Inventário de múltiplas etiquetas e anticolisão}

Quando várias tags recebem uma consulta, respostas simultâneas podem colidir. O EPC Gen2 organiza o inventário por um protocolo probabilístico de slots. O interrogador envia uma consulta contendo o parâmetro $Q$; cada tag participante escolhe aleatoriamente um contador entre 0 e $2^Q-1$. Tags cujo contador é zero respondem, enquanto as demais aguardam comandos que avançam ou ajustam a rodada \cite{gs1gen2}.

Assim, o anticolisão da interface aérea é executado pelo firmware do leitor em cooperação com as tags. Não é correto descrever o Arduino como responsável por resolver colisões de RF. As responsabilidades do programa embarcado são:

\begin{itemize}
    \item iniciar e interromper a rodada de inventário;
    \item configurar, quando suportado, parâmetros de consulta e sessão;
    \item receber todos os quadros produzidos pelo módulo;
    \item validar o enquadramento e o byte de verificação;
    \item eliminar EPCs duplicados dentro da janela de inventário;
    \item registrar contagem, tempo, RSSI bruto e erros de recepção.
\end{itemize}

Algoritmos acadêmicos podem melhorar a escolha de parâmetros ou substituir estratégias de inventário em leitores experimentais \cite{su2019anticollision}. No protótipo deste trabalho, contudo, o ponto de partida deve ser o mecanismo Gen2 já implementado no YPD-R200. A contribuição embarcada está na aquisição confiável e na avaliação de sua capacidade, não na reimplementação da camada física.

\section{Antenas e formação da zona de leitura}

\subsection{Ganho e largura de feixe}

O ganho em dBi expressa a concentração da radiação em determinadas direções em relação a uma fonte isotrópica. Em geral, maior ganho concentra mais energia no lóbulo principal, podendo ampliar o alcance no eixo de máxima radiação em troca de menor cobertura angular \cite{rao2005uhf}. Para inventariar uma sala, uma zona mais larga pode ser mais útil que a maior distância possível; para um ponto fixo em uma porta, um feixe controlado reduz leituras de áreas adjacentes.

Ganho, polarização e cobertura espacial são propriedades diferentes. Polarização linear ou circular descreve a trajetória do campo elétrico; cobertura direcional ou omnidirecional descreve como a potência é distribuída no espaço. Portanto, ``linear'' não é sinônimo de ``direcional'', assim como ``omnidirecional'' não é o oposto de ``linear''. Uma antena omnidirecional pode, inclusive, possuir polarização linear.

A antena especificada na arquitetura e adotada nos ensaios é a Airplux APCA8090. A página do fabricante lista as frequências centrais de 915~MHz e 866,5~MHz, ganho de 5~dBi, polarização linear, razão axial menor ou igual a 5~dB, impedância de 50~$\Omega$, conector SMA ou customizado, elemento cerâmico de 80~$\times$~80~$\times$~7~mm e placa de 90~$\times$~90~mm \cite{airplux2026apca8090}. A documentação consultada não torna inequívoco se 915~MHz e 866,5~MHz correspondem a variantes distintas, nem informa largura de banda, VSWR, largura de feixe ou relação entre frente e costas. Assim, esses valores constituem referência nominal, não comprovação da resposta da unidade utilizada em toda a faixa de operação.

\subsection{Polarização e orientação}

Em uma antena de polarização linear, o vetor de campo elétrico mantém uma orientação predominante. Quando a antena da etiqueta também é linear, o desalinhamento angular entre as polarizações reduz o acoplamento e pode diminuir a probabilidade de leitura. Por isso, a orientação da AD-239 M730 em relação à APCA8090 foi incluída como variável experimental e aparece nos ensaios de orientação do conjunto.

Como a APCA8090 é especificada como antena de polarização linear, a orientação relativa entre antena do leitor e antena da etiqueta é uma variável potencialmente crítica. Os ensaios de $0^\circ$, $45^\circ$ e $90^\circ$ avaliam essa dependência nas condições específicas da campanha.

Quando a orientação das etiquetas não é controlada, a intervenção mais diretamente relacionada ao problema é empregar polarização circular ou diversidade de polarização \cite{times72022polarization}. Um padrão omnidirecional trata outra dimensão do projeto: pode ampliar a cobertura no plano azimutal e reduzir a necessidade de apontamento, mas não elimina por si só a incompatibilidade de polarização. Como antenas não foram comparadas nesta campanha, uma eventual vantagem omnidirecional para a varredura portátil permanece hipótese a testar. Maior cobertura ao redor do leitor também pode elevar leituras externas e a exposição ao multipercurso; para pontos fixos, o controle espacial de uma antena direcional pode continuar preferível.

\subsection{Multipercurso e diversidade}

Em ambientes fechados, o sinal transmitido pelo leitor não se propaga somente pelo caminho direto entre a antena e a etiqueta. Parte da energia é refletida por paredes, piso, teto, móveis, equipamentos e estruturas metálicas, formando componentes que percorrem distâncias diferentes e chegam à tag com amplitudes e fases próprias. A superposição pode ser construtiva, elevando o campo resultante, ou destrutiva, reduzindo-o em pontos específicos da zona de leitura.

Tomando a fase do caminho direto como referência, uma representação fasorial escalar, complexa e de banda estreita da componente do campo alinhada à polarização da tag é

\begin{equation}
    \widetilde{E}_{\mathrm{tot}}
    = E_{0} + \sum_{i=1}^{N} E_i e^{-j\phi_i},
    \label{eq:campo-multipercurso}
\end{equation}

em que $E_0$ representa a amplitude do caminho dominante, $E_i$ é a amplitude da $i$-ésima componente refletida e $\phi_i$ é sua fase relativa. A potência disponível na antena da etiqueta depende, entre outros fatores, de $\lvert\widetilde{E}_{\mathrm{tot}}\rvert^2$, da polarização e do casamento de impedâncias. Assim, componentes aproximadamente em oposição de fase podem produzir um nulo local e reduzir a potência abaixo do limiar necessário para energizar o chip, mesmo que a tag esteja dentro de uma distância que seria considerada adequada por um modelo de espaço livre. Em RFID UHF passivo, essa representação ilustra principalmente o enlace de ida, responsável por energizar a tag; a leitura depende também do enlace de retorno por retroespalhamento, igualmente sujeito à propagação. A \autoref{eq:campo-multipercurso} é apenas conceitual: uma análise eletromagnética completa exigiria campos vetoriais, coeficientes de reflexão, geometria e propriedades dos materiais.

Quando existe uma componente dominante acompanhada por muitas componentes espalhadas, o desvanecimento de pequena escala pode ser representado conceitualmente por um canal Rician. Com potência média normalizada,

\begin{equation}
    h = \sqrt{\frac{K}{K+1}}\,e^{j\phi_0}
        + \sqrt{\frac{1}{K+1}}\,g,
    \qquad g \sim \mathcal{CN}(0,1),
    \label{eq:canal-rician}
\end{equation}

em que a primeira parcela representa a componente dominante, $g$ é uma variável gaussiana complexa normalizada que representa a parcela difusa e $K$ é a razão linear entre a potência dominante e a potência espalhada. Um valor de $K$ expresso em decibéis precisaria ser convertido antes de ser usado na \autoref{eq:canal-rician}. Na ausência de componente dominante, o caso limite $K=0$ conduz ao modelo Rayleigh para a envoltória, desde que sejam válidas as hipóteses estatísticas de espalhamento. Esses modelos descrevem distribuições possíveis do canal; não constituem, por si só, prova de que o laboratório deste trabalho apresentou uma distribuição Rician ou Rayleigh.

A fase acumulada em um percurso depende de seu comprimento e da frequência. Para o $i$-ésimo caminho, pode-se escrever

\begin{equation}
    \phi_i = \frac{2\pi(d_i-d_0)}{\lambda} + \theta_i
           = \frac{2\pi f(d_i-d_0)}{c} + \theta_i
           \pmod{2\pi},
    \label{eq:fase-percurso}
\end{equation}

em que $d_i$ é o comprimento do percurso refletido, $d_0$ é o comprimento do caminho dominante, $\lambda=c/f$ é o comprimento de onda e $\theta_i$ reúne as mudanças adicionais de fase introduzidas pelas reflexões e pelas antenas. Alterar a frequência modifica as fases relativas dos caminhos; deslocar a antena ou a tag modifica seus comprimentos relativos; e utilizar antenas em posições distintas modifica a distribuição espacial dos nulos. Em uma arquitetura com múltiplas antenas, também é possível variar a fase relativa entre os caminhos de alimentação. Uma mudança de fase comum aplicada a uma única antena, entretanto, não altera por si só a interferência relativa entre os percursos.

\citeonline{sabesan2014diversity} modelaram o desvanecimento Rician e adotaram o multipercurso como motivação para um sistema RFID UHF passivo de ampla cobertura que combina duas ou mais antenas espacialmente separadas com saltos controlados de fase e de frequência. A finalidade da diversidade é oferecer configurações de canal distintas: uma tag situada em um nulo em determinada frequência, posição de antena ou relação de fases pode receber potência suficiente em outra configuração. Isso reduz a dependência de um único estado do canal, mas não garante a eliminação de todos os nulos. A arquitetura daquele estudo empregou um transceptor Impinj R2000, que não deve ser confundido com o módulo YPD-R200 utilizado neste trabalho.

No protótipo deste trabalho, essa fundamentação justifica padronizar e registrar a geometria e avaliar repetibilidade em vez de uma leitura isolada. A altura foi mantida em 150~cm nos ensaios de distância, orientação e material, mas não foi variada como fator experimental. O YPD-R200 adquirido possui uma única conexão SMA; portanto, a diversidade espacial completa com múltiplas antenas exigiria chaveamento externo, mais de um módulo ou outra classe de leitor. A campanha não mediu resposta impulsiva, fase, fator $K$ ou coeficientes de reflexão, e o código do Arduino não estimou nem compensou o canal. A região, os canais efetivos e o FHSS também não foram confirmados durante a coleta; portanto, eventual salto de frequência do módulo não é apresentado como implementação da técnica de \citeonline{sabesan2014diversity}. As repetições e o deslocamento do leitor no inventário permitiram observar o desempenho agregado em um cenário institucional realista. Nesse contexto, o multipercurso integra as condições normais de operação; ele permanece uma explicação concorrente para a variabilidade, e não uma variável quantificada separadamente.

\section{Etiqueta AD-239 M730}

Uma etiqueta completa contém o circuito integrado e uma antena, formando o \textit{inlay}; o conjunto pode ainda receber adesivo, substrato e acabamento de etiqueta. Informar apenas o chip é insuficiente para prever alcance, pois o desempenho prático depende também da antena do inlay, da superfície de instalação e da orientação em relação à antena do leitor \cite{rao2005uhf}.

A AD-239 M730, adotada como etiqueta de referência do trabalho, utiliza o chip Impinj M730, opera em UHF entre 860 e 960~MHz, é compatível com ISO/IEC 18000-63 e EPC Class 1 Gen2, possui 128 bits de memória EPC, TID de 96 bits e número serial único de 48 bits embutido no TID \cite{avery2022ad239m730}. Sua antena possui 70~$\times$~14,5~mm e o produto é fornecido em formatos como \textit{dry inlay}, \textit{wet inlay} e etiqueta.

O datasheet classifica a AD-239 M730 como solução para superfícies não metálicas. Por esse motivo, o papelão foi usado como condição de referência no ensaio específico de material, enquanto a aplicação diretamente sobre metal foi tratada como condição crítica de degradação. Uma terceira condição separou a etiqueta do metal por um espaçador plástico de 10~mm. Nos ensaios de distância e orientação, papelão, madeira e plástico permaneceram associados, respectivamente, a TAG1, TAG2 e TAG3; por isso, esses suportes não foram avaliados como fatores independentes nesses dois blocos. Vidro e outras espessuras não integraram a análise experimental. A \autoref{tab:ad239m730} reúne as especificações relevantes para esse recorte.

\begin{table}[htbp]
    \centering
    \caption{Especificações relevantes da etiqueta AD-239 M730.}
    \label{tab:ad239m730}
    \footnotesize
    \begin{tabularx}{\textwidth}{p{4.0cm}X}
        \toprule
        \textbf{Parâmetro} & \textbf{Especificação da AD-239 M730} \\
        \midrule
        Faixa de frequência & UHF de 860 a 960~MHz. \\
        Circuito integrado & Impinj M730. \\
        Padrão de operação & ISO/IEC 18000-63 e EPC Class 1 Gen2. \\
        Memória EPC & 128 bits. \\
        Memória TID & 96 bits, com número serial único de 48 bits embutido no TID. \\
        Dimensão da antena do inlay & 70~$\times$~14,5~mm. \\
        Formatos de fornecimento & \textit{Dry inlay}, \textit{wet inlay} e etiqueta. \\
        Aplicação indicada & Bens não metálicos ou superfícies de baixa influência condutiva. \\
        Limitação considerada & Aplicação direta sobre metal pode reduzir ou impedir a leitura, devendo ser avaliada apenas como condição crítica. \\
        \bottomrule
    \end{tabularx}
    \fonte{\citeonline{avery2022ad239m730}.}
\end{table}

Para bens metálicos, uma etiqueta convencional pode ser dessintonizada pela superfície condutora. Como a AD-239 M730 é especificada para uso não metálico, ela não deve ser apresentada como solução definitiva para computadores, armários, instrumentos com carcaça metálica ou outros bens condutivos. Neste trabalho, o metal entra como condição crítica de degradação observada no ensaio.

\section{Módulo YPD-R200}

O YPD-R200 é um módulo embarcado compatível com ISO 18000-6C/EPC Class 1 Gen2. O datasheet do fabricante informa alimentação entre 3,6 e 5,5~V, UART de até 115200~bit/s, dimensões de 53~$\times$~33~mm e versões com potência máxima de 20 ou 26~dBm \cite{yanpodo2021r200}. A variante empregada possui potência máxima nominal de 26~dBm, equivalente a aproximadamente 398~mW de RF no conector do módulo, e corrente de pico informada em aproximadamente 380~mA. Com uma antena de ganho nominal de 5~dBi, a EIRP teórica seria de até 31~dBm antes de descontar as perdas de cabo e conectores. Esse valor é uma estimativa de orçamento de enlace, não uma potência irradiada medida nem uma comprovação de conformidade regulatória.

Para essa variante, o fabricante apresenta alcance típico de 3~m em espaço aberto com antena cerâmica de 45~$\times$~45~mm. Esse valor é mais defensável que alegações genéricas de 15 ou 20~m encontradas em páginas de terceiros. A APCA8090 possui ganho nominal maior e dimensões superiores às da antena usada no exemplo do datasheet, mas os ensaios deste trabalho não estabelecem o alcance real absoluto; eles apenas mostram o comportamento do conjunto nas distâncias testadas.

O protocolo serial utiliza quadros delimitados por \texttt{0xAA} e \texttt{0xDD}. Entre eles aparecem tipo, comando, tamanho dos parâmetros, parâmetros e soma de verificação. O comando \texttt{0x22} executa uma consulta de inventário, \texttt{0x27} inicia múltiplas consultas e \texttt{0x28} interrompe essa operação. Além dos comandos de inventário, o comando \texttt{0xB6} permite configurar a potência de transmissão, e as notificações retornadas incluem RSSI, PC, EPC e CRC da tag \cite{yanpodo2016protocol}.

\section{Arquiteturas de aplicação}

\subsection{Inventário portátil}

No cenário principal, leitor, antena, microcontrolador e fonte formam uma unidade transportada pelo operador. A cada ambiente é aberta uma janela de inventário; os EPCs únicos são acumulados e comparados com a relação esperada. A ergonomia do equipamento, a alimentação e a estabilidade da UART são requisitos tão importantes quanto a distância nominal.

O Arduino Uno/Nano permite validar a comunicação, mas sua única UART de hardware é compartilhada com a USB. O uso de \texttt{SoftwareSerial} a 115200~bit/s aumenta o risco de perda de bytes. Para uma versão portátil confiável, Arduino Mega, ESP32 ou outra plataforma com UART adicional em hardware é tecnicamente mais adequada.

\subsection{Pontos fixos de movimentação}

Um leitor fixo pode registrar que uma tag esteve em sua zona de cobertura, criando um evento de passagem. Para transformar esse evento em ``entrada'' ou ``saída'', é necessário distinguir a ordem espacial das observações. Pesquisas de portais RFID empregam duas antenas, arranjos direcionais, sensores auxiliares ou séries temporais de RSSI e fase \cite{wang2022rfaccess}.

Consequentemente, o YPD-R200 com uma única antena pode servir como unidade de detecção, mas não constitui sozinho um controlador de direção. Uma arquitetura futura pode usar duas zonas consecutivas, cada uma com timestamp, e inferir o sentido pela ordem de ativação. Essa extensão também requer tratamento de leituras repetidas e de tags localizadas perto da porta, mas que não a atravessaram.
